% Einheit 2 von 5 – Extrempunkte (bank/kurvenuntersuchung/e2.jsonl, zusammenbau v0.1)
%% TODO Zeitmarke Einheit 2: Marken-Zeile nicht lesbar
%% TODO Zweigzeile Teil 1: was hier gelernt wird (Satz in Schülersprache aus dem Einheitstitel) – Einheit 2
\einheitenkopf[e2]{Einheit 2 von 5 · Extrempunkte}
\zweigzeile{Abitur GK}
\setcounter{aufgabe}{10}

%% TODO \verfahren{…}: Verfahrensüberschrift in Sachsprache fehlt in der Bank (2.3 b)
%% TODO Titel: Ich-kann-Satz fehlt in der Bank (Platzhalter: Kettenname) – Nr. 11
%% TODO Anweisung: ein Satz über der ersten Teilaufgabe fehlt in der Bank – Nr. 11
\begin{aufgabe}{Extrempunkte berechnen}
%% TODO \rechenplatz{n}: Schrittzahl des Grundfalls fehlt in der Bank (nur bei fester Schreibform, 2.3 b)
\begin{teile}
\teil Kreuze an, ohne zu rechnen: „Zeigen Sie, dass der Graph bei $x = 2$ einen Hochpunkt hat.“ \\ \kreuz{Stelle gegeben – nachweisen} \\ \kreuz{Stelle gesucht – berechnen}
\end{teile}
\begin{gleichungsraster}[2]
\gl{\text{Berechne die Stellen, an denen der Graph von $f(x) = x^3 - 1{,}5x^2 - 6x + 2$ eine waagerechte Tangente hat.}} &
\end{gleichungsraster}
\begin{teile}
\teil Bestimme Lage und Art der Extrempunkte des Graphen von $f(x) = x^3 + 1{,}5x^2 - 6x + 1$.
\teil Bestimme Lage und Art aller Extrempunkte des Graphen von $f(x) = x^4 - 8x^2$.
\teil Bestimme Koordinaten und Art der beiden Extrempunkte des Graphen von $f(x) = 3x^5 - 20x^3$ und weise den Sattelpunkt nach. \hfill (FHR 2025)
\teil Der Graph von $f(x) = x^4 - 4x^3 - 2x^2 + 12x$ hat bei $x = 3$ eine Extremstelle. Bestimme die weiteren Extremstellen und Koordinaten und Art aller Extrempunkte. \hfill (FHR 2023)
\teil Bestimme die Koordinaten des Hochpunkts des Graphen von $f(x) = (3 - x) \cdot \mathrm{e}^{x}$. \hfill (Abitur 2025 GK)
\teil Die Abbildung zeigt den Graphen von $f(x) = x \cdot \mathrm{e}^{-0{,}5x}$. Berechne die Koordinaten seines Extrempunkts; auf die hinreichende Bedingung wird verzichtet. Begründe die Art mit der Abbildung.

\begin{ksys}[xmin=0,xmax=6,ymin=-1,ymax=1,ablesen] \funktion{\x*exp(-0.5*\x)}{f} \end{ksys}
\teil Die Randlinie eines liegenden Kerzenhalters wird für $0 \le x \le 4$ durch $r(x) = 0{,}25x^3 - 1{,}5x^2 + 2{,}25x + 1$ beschrieben ($x$ und $r(x)$ in cm, $r(x)$ ist der Radius an der Stelle $x$). Ermittle rechnerisch die Stellen, an denen der Radius am größten ist. \hfill (Abitur 2018 LK)
\teil Gib den Wertebereich von $f(x) = x^4 - 2x^2$ an. Bestimme dazu die Tiefpunkte und das Verhalten für $x \to \pm\infty$.
\teil Gegeben ist $f(x) = (x - 2)^2 \cdot \mathrm{e}^{-0{,}5x}$; ohne Nachweis gilt $f''(x) = (0{,}25x^2 - 3x + 7) \cdot \mathrm{e}^{-0{,}5x}$. Bestimme Koordinaten und Art der lokalen Extrempunkte des Graphen. \hfill (Abitur 2021 GK)
\end{teile}
\end{aufgabe}

%% TODO \verfahren{…}: Verfahrensüberschrift in Sachsprache fehlt in der Bank (2.3 b)
%% TODO Titel: Ich-kann-Satz fehlt in der Bank (Platzhalter: Kettenname) – Nr. 12
%% TODO Anweisung: ein Satz über der ersten Teilaufgabe fehlt in der Bank – Nr. 12
\begin{aufgabe}{Extrempunkte nachweisen}
%% TODO \rechenplatz{n}: Schrittzahl des Grundfalls fehlt in der Bank (nur bei fester Schreibform, 2.3 b)
\begin{teile}
\teil Kreuze an, ohne zu rechnen: „Weisen Sie nach, dass der Graph bei $x = -1$ einen Tiefpunkt hat.“ \\ \kreuz{Stelle gegeben – nachweisen} \\ \kreuz{Stelle gesucht – berechnen}
\teil Zeige, dass $f'(2) = 0$ gilt, der Graph von $f(x) = 2x^3 - 6x^2$ also bei $x = 2$ eine waagerechte Tangente hat.
\teil Zeige, dass der Graph von $f(x) = x^3 - 27x + 1$ einen Hochpunkt mit der $x$-Koordinate $-3$ hat. \hfill (Abitur 2026 GK)
\teil Zeige, dass $H(2 | 5)$ ein Hochpunkt des Graphen von $f(x) = -x^3 + 3x^2 + 1$ ist.
\teil Weise nach, dass $0$ und $2$ die Extremstellen von $f(x) = x^3 - 3x^2 + 2$ sind. Nutze den Vorzeichenwechsel von $f'$. \hfill (Abitur 2020 GK)
\teil Gegeben ist $f(x) = \frac{1}{3}x^3 - 9x$; es gilt $f''(3) \ne 0$. Weise nach, dass $3$ eine Extremstelle von $f$ ist. \hfill (Abitur 2025 LK)
\teil Zeige, dass der Graph von $f(x) = x^3 - 6x^2 + 12x$ bei $x = 2$ eine waagerechte Tangente, aber keinen Extrempunkt hat. Begründe ohne weitere Rechnung mit der Form von $f'$.
\teil Begründe, dass der Graph von $f(x) = x^4 + 3x^2 - 2x$ genau einen Tiefpunkt hat.
\teil Zeige, dass $f(x) = x^3 - 5x$ zwischen $1$ und $2$ eine Extremstelle hat, ohne sie zu berechnen.
\teil Begründe, dass der Graph von $f(x) = \mathrm{ln}(x^2 + 4)$ bei $x = 0$ einen Tiefpunkt hat, und gib ihn an.
\end{teile}
\end{aufgabe}

%% TODO Titel: Ich-kann-Satz fehlt in der Bank (Platzhalter: Kettenname) – Nr. 13
%% TODO Anweisung: ein Satz über der ersten Teilaufgabe fehlt in der Bank – Nr. 13
\begin{aufgabe}{Abstand zweier Extrempunkte verschiedener Graphen mit einer Schranke vergleichen}
\begin{teile}
\teil Gegeben sind $f(x) = -x^3 + 3x^2 + 2$ mit dem Hochpunkt $H(2 | 6)$ und die Parabel $p(x) = -x^2 + 6x - 3{,}5$, deren Scheitel ein Hochpunkt ist. Zeige, dass der Abstand der beiden Hochpunkte größer als $0{,}5$ ist. \hfill (Abitur 2021 GK)
\end{teile}
\end{aufgabe}

%% TODO Titel: Ich-kann-Satz fehlt in der Bank (Platzhalter: Kettenname) – Nr. 14
%% TODO Anweisung: ein Satz über der ersten Teilaufgabe fehlt in der Bank – Nr. 14
\begin{aufgabe}{Abstand zweier Extrempunkte über die Punktsymmetrie berechnen}
\begin{teile}
\teil Der Graph von $f(x) = 2x^3 - 6x$ ist symmetrisch zum Ursprung und hat einen Hochpunkt mit der $x$-Koordinate $-1$. Berechne den Abstand zwischen Hoch- und Tiefpunkt. \hfill (Abitur 2026 GK)
\end{teile}
\end{aufgabe}

%% TODO Titel: Ich-kann-Satz fehlt in der Bank (Platzhalter: Kettenname) – Nr. 15
%% TODO Anweisung: ein Satz über der ersten Teilaufgabe fehlt in der Bank – Nr. 15
\begin{aufgabe}{Extrempunkte berechnen und ihre Verbindungsgerade als Winkelhalbierende nachweisen}
\begin{teile}
\teil Der Graph von $f(x) = \frac{1}{8}x^3 - 1{,}5x$ hat genau zwei Extrempunkte. Berechne ihre Koordinaten und weise nach, dass die Gerade durch sie die Winkelhalbierende des II. und IV. Quadranten ist. \hfill (Abitur 2025 GK)
\end{teile}
\end{aufgabe}

%% TODO Titel: Ich-kann-Satz fehlt in der Bank (Platzhalter: Kettenname) – Nr. 16
%% TODO Anweisung: ein Satz über der ersten Teilaufgabe fehlt in der Bank – Nr. 16
\begin{aufgabe}{Nullstellen und Extremstelle einer Parabel berechnen}
\begin{teile}
\teil Gegeben ist $r(x) = -(x^2 - 2x - 3)$. Berechne die Nullstellen und die Extremstelle von $r$. \hfill (Abitur 2023 GK)
\end{teile}
\end{aufgabe}

%% TODO Titel: Ich-kann-Satz fehlt in der Bank (Platzhalter: Kettenname) – Nr. 17
%% TODO Anweisung: ein Satz über der ersten Teilaufgabe fehlt in der Bank – Nr. 17
\begin{aufgabe}{Berührung des Graphen mit der x-Achse über die Extrempunkte begründen}
\begin{teile}
\teil Begründe mit den Extrempunkten, dass der Graph von $f(x) = x^3 - 3x^2 + 4$ die $x$-Achse berührt.
\end{teile}
\end{aufgabe}

%% TODO Titel: Ich-kann-Satz fehlt in der Bank (Platzhalter: Kettenname) – Nr. 18
%% TODO Anweisung: ein Satz über der ersten Teilaufgabe fehlt in der Bank – Nr. 18
\begin{aufgabe}{Größten Funktionswert über Monotonie, Grenzverhalten und Symmetrie begründen und Wertemenge angeben}
\begin{teile}
\teil Begründe, dass $f(x) = 5\,\mathrm{e}^{-0{,}5x^2}$ bei $x = 0$ seinen größten Wert hat, und gib die Wertemenge an.
\end{teile}
\end{aufgabe}

%% TODO Titel: Ich-kann-Satz fehlt in der Bank (Platzhalter: Kettenname) – Nr. 19
%% TODO Anweisung: ein Satz über der ersten Teilaufgabe fehlt in der Bank – Nr. 19
\begin{aufgabe}{Extremstelle zwischen zwei Stellen mit gleichem Funktionswert ohne Rechnung begründen}
\begin{teile}
\teil Für eine ganzrationale, nicht konstante Funktion $f$ gilt $f(2) = f(5)$. Begründe ohne Rechnung, dass $f$ zwischen $2$ und $5$ mindestens eine Extremstelle hat. \hfill (Abitur 2018 LK)
\end{teile}
\end{aufgabe}

%% TODO Titel: Ich-kann-Satz fehlt in der Bank (Platzhalter: Kettenname) – Nr. 20
%% TODO Anweisung: ein Satz über der ersten Teilaufgabe fehlt in der Bank – Nr. 20
\begin{aufgabe}{Lage eines Graphen zwischen zwei Geraden auf einem Intervall über die Differenzfunktionen nachweisen}
\begin{teile}
\teil Gegeben sind $f(x) = -0{,}25x^3 + 1{,}5x^2 - 1{,}75x + 1$ und die Geraden $g_o(x) = 0{,}5x + 1$ und $g_u(x) = 0{,}5x - 2$. Weise nach, dass der Graph von $f$ für $0 \le x \le 4$ zwischen den beiden Geraden liegt, und bestimme die Stelle im Inneren mit dem kleinsten senkrechten Abstand zu $g_u$ und diesen Abstand. \hfill (Abitur 2017 GK)
\end{teile}
\end{aufgabe}

%% TODO Titel: Ich-kann-Satz fehlt in der Bank (Platzhalter: Kettenname) – Nr. 21
%% TODO Anweisung: ein Satz über der ersten Teilaufgabe fehlt in der Bank – Nr. 21
\begin{aufgabe}{Aussagen über Funktions- und Ableitungswert nahe dem Tiefpunkt ohne Rechnung beurteilen}
\begin{teile}
\teil Der Graph einer ganzrationalen Funktion $f$ hat den Tiefpunkt $T(2 | -3)$ und keinen weiteren Extrempunkt. Beurteile ohne Rechnung: „$f'(1{,}9) > 0$“ \\ \kreuz{wahr} \\ \kreuz{falsch}
\end{teile}
\end{aufgabe}

%% TODO Titel: Ich-kann-Satz fehlt in der Bank (Platzhalter: Kettenname) – Nr. 22
%% TODO Anweisung: ein Satz über der ersten Teilaufgabe fehlt in der Bank – Nr. 22
\begin{aufgabe}{Extrempunkte berechnen – Fehler finden, Begründen, Anwendung}
\begin{teile}
\teil Finde den Fehler und stelle richtig. Gesucht war die Art des Extrempunkts von $f(x) = x^3 - 12x + 1$ bei $x = 2$. \rechnung{f''(x) &= 6x \\ f''(2) &= 12 > 0 &&\Rightarrow \text{Hochpunkt}}
\teil Begründe, warum $f'(x_0) = 0$ allein noch nicht zeigt, dass bei $x_0$ ein Extrempunkt liegt.
\teil Die Decke eines Tunnels wird im Querschnitt durch $p(x) = -0{,}25x^2 + 2x + 1$ beschrieben ($x$ und $p(x)$ in m). Ein Lkw ist $4$ m hoch. Beschreibe, wie man den Abstand zwischen Lkw-Dach und Tunneldecke an der höchsten Stelle ermittelt, und berechne ihn. \hfill (Abitur 2023 GK)
\end{teile}
\end{aufgabe}
